% Section 25.3: Geodesics, Killing fields, Noether and the inertia metric | Part: 5 | Depends on: part:manifolds, part:integration, sec:lc | Status: done | Original lines: 3404-3422
%!TEX root = ../../../main.tex
\subsection{Geodesics, Killing fields, Noether and the inertia metric}\label{sec:noethersynth}

The Lagrangian mechanics of Section~\ref{sec:phys-energy} and the geodesic equations of Section~\ref{sec:lc} are the same equations.

\begin{synthesis}[free motion $=$ geodesic of the inertia metric]
For $L=\frac12g(v,v)-V$ with $V=0$, the Euler--Lagrange equations are the geodesic equations of $g$ (Section~\ref{sec:lc}, via the energy functional). The Christoffel symbols are the ``fictitious forces'' of non-inertial coordinates.
\end{synthesis}

\begin{example}[the free particle in polar coordinates, both ways]\label{ex:polarfree}
% REWRITTEN: the geometry side is the computation of Examples ex:polarstraight, ex:polarEL and ex:angmom (Section sec:conn, sec:lc), with \varphi written \theta here as in the rest of Part V.
Mechanics: $T=\frac12m(\dot r^2+r^2\dot\theta^2)$ gives $\frac{\dd}{\dd t}(m\dot r)-mr\dot\theta^2=0$ and $\frac{\dd}{\dd t}(mr^2\dot\theta)=0$.
Geometry: for $g=\dd r^2+r^2\dd\theta^2$ the geodesic equations were found in Example~\ref{ex:polarstraight} (and from the Lagrangian in Example~\ref{ex:polarEL}): $\ddot r-r\dot\theta^2=0$ and $\frac{\dd}{\dd t}(r^2\dot\theta)=0$.
The two systems coincide; $-r\dot\theta^2$ is the centrifugal term. The vector field $\partial_\theta$ is Killing ($\Lie_{\partial_\theta}g=0$), and the conserved quantity of a geodesic, $g(\partial_\theta,\dot\gamma)=r^2\dot\theta$, is the angular momentum $p_\theta=p(\partial_\theta)$ of Section~\ref{sec:phys-energy} (Example~\ref{ex:angmom}).
\end{example}

\begin{keyidea}[one conservation law, three names]
\emph{Mechanics:} if the flow of $\xi$ leaves $L$ invariant, $J_\xi=p(\xi)$ is conserved (Noether). \emph{Geometry:} if $\xi$ is Killing, $g(\xi,\dot\gamma)$ is constant along geodesics. \emph{Relativity:} $E=-p(\partial_t)$ for a stationary spacetime.
They are the same statement because $p=g(\dot\gamma,\cdot)$ for the kinetic Lagrangian; the metric is the inertia of Section~\ref{sec:phys-energy}, and Killing fields are the symmetries of the Lagrangian.
\end{keyidea}
